\documentclass[a4paper]{article}
\usepackage{amsfonts}
\usepackage{amsmath}
\usepackage{amssymb}
\usepackage{graphicx}     
\begin{document}
	
	\noindent \large Auteur: Ilyas Sal-lak \\
	
	\noindent \large Studierichting: Informatica\\
	
	\noindent \large Oefeningenreeks 1D nr. 46.2\\
	
	\medskip
	\normalsize
	Voor $z^4 = -1$ zoeken we de vierdemachtswortels.\\
	
	Formule: $z = \left(\cos\left(\dfrac{\theta + 2\pi k}{4}\right) + i\sin\left(\dfrac{\theta + 2\pi k}{4}\right)\right)$\\
	
	In poolcoördinaten geldt $r = 1$ en $\theta = \pi$\\
	
	Voor $k = 0, 1, 2, 3$ vinden we:\\
	
	$z_1 = \left(\cos\left(\dfrac{\pi}{4}\right) + i\sin\left(\dfrac{\pi}{4}\right)\right) = \dfrac{\sqrt{2}}{2} + i\dfrac{\sqrt{2}}{2}$\\
	
	$z_2 = \left(\cos\left(\dfrac{3\pi}{4}\right) + i\sin\left(\dfrac{3\pi}{4}\right)\right) = -\dfrac{\sqrt{2}}{2} + i\dfrac{\sqrt{2}}{2}$\\
	
	$z_3 = \left(\cos\left(\dfrac{5\pi}{4}\right) + i\sin\left(\dfrac{5\pi}{4}\right)\right) = -\dfrac{\sqrt{2}}{2} - i\dfrac{\sqrt{2}}{2}$\\
	
	$z_4 = \left(\cos\left(\dfrac{7\pi}{4}\right) + i\sin\left(\dfrac{7\pi}{4}\right)\right) = \dfrac{\sqrt{2}}{2} - i\dfrac{\sqrt{2}}{2}$\\
	
	De vier vierdemachtswortels zijn dus:\\
	\begin{align*}
		z_1 &= \dfrac{\sqrt{2}}{2} + i\dfrac{\sqrt{2}}{2} \\
		z_2 &= -\dfrac{\sqrt{2}}{2} + i\dfrac{\sqrt{2}}{2} \\
		z_3 &= -\dfrac{\sqrt{2}}{2} - i\dfrac{\sqrt{2}}{2} \\
		z_4 &= \dfrac{\sqrt{2}}{2} - i\dfrac{\sqrt{2}}{2} \\
	\end{align*}
	
	\begin{thebibliography}{999}
		%\bibitem{a} Referentie
	\end{thebibliography}
\end{document}